\documentclass[10pt,a4paper]{amsart}
\usepackage[margin=19mm]{geometry}
\usepackage{amsmath,amssymb,mathtools}
\usepackage[scr=rsfs,cal=euler]{mathalfa}
\usepackage{xfrac}
\usepackage[unicode,hidelinks,bookmarks=false]{hyperref}
\newcommand{\ie}{i.e.\ }
\newcommand{\cf}{cf.\ }
\newcommand{\resp}{respectively\ }
\newcommand{\Subsection}{\subsection}
\providecommand{\nobreakdash}{\nobreak\hbox{-}}
\setlength{\emergencystretch}{2em}
\multlinegap0pt
\usepackage{amssymb}
\usepackage{mathtools}
\usepackage[shortlabels]{enumitem}
\usepackage{xcolor}
\newtheorem{proposition}{Proposition}
\newcommand{\Om}{\Omega}
\newcommand{\PhiP}{\Phi_{p,\beta}}
\newcommand{\SigmaP}{\Sigma_p}
\newcommand{\dd}{\,\mathrm d}
\title{Non-Shattering at and Above the Dynamical Temperature in the Spherical Pure \(p\)-Spin Model}
\author{Taegyun Kim}
\date{Source: arXiv:2608.14369v2 (CC BY 4.0)}
\begin{document}
\maketitle

\noindent\emph{Source and licence.} Non-Shattering at and Above the Dynamical Temperature in the Spherical Pure \(p\)-Spin Model. Version arXiv:2608.14369v2 (CC BY 4.0), distributed by arXiv under the Creative Commons Attribution 4.0 International licence, http://creativecommons.org/licenses/by/4.0/, as recorded in the arXiv licence metadata for this exact version and shown on the abstract page https://arxiv.org/abs/2608.14369v2. Full licence notice: \texttt{licenses/arxiv-2608.14369-CC-BY-4.0.txt}.\par\smallskip

\noindent\textit{Editorial scope.} Counted unit: the article's proposition prop:general-p-marked (source0.tex lines 1076--1094). The article's complete original proof of it is delivered with it (source0.tex lines 1096--1261, one \texttt{proof} environment of 166 lines). No mathematics is written, repaired or re-derived: the statement and the proof are the article's own lines, in the article's own notation, and no retained line is substituted or re-flowed.  Retained uncounted as the source-local context the proof needs: lines 691--694; lines 708--713.  Nothing was omitted from any retained span, so the package carries no omission record.  No retained line contains a citation, so the wrapper carries no bibliography and no bibliography entry.  No retained line contains a figure environment, an image-inclusion command or drawing code, so no image is delivered and the package declares no asset dependency.  Editorial formatting only, in addition: an AMS \texttt{amsart} preamble that restates the article's own declarations and macros used by the retained lines, namely the environments \texttt{proposition} on separate counters, together with the standard packages those lines need.  The retained environments print in this document as Proposition 1 in order of appearance, whereas the article prints the same items with the numbering of its own full document.  The complete native source of the article as distributed in the arXiv e-print for this exact version is bundled unchanged as \texttt{sources/207665/source0.tex}.
\begin{equation}\label{eq:Omega}
 \Om(x)=\int_{-2}^2\log\lvert x-\lambda\rvert\,
 \mu_{\mathrm{sc}}(\dd\lambda).
\end{equation}

\begin{equation}\label{eq:PhiGeneral}
 \PhiP(u,t)
 =\SigmaP(u)+\beta t^p u
  +\frac12\log(1-t^2)
  +\frac{\beta^2}{2}v_p(t).
\end{equation}

\begin{proposition}[General-\(p\) marked sign law]
\label{prop:general-p-marked}
Fix \(p\geq3\) and \(0<\beta\leq\beta_{\mathrm{sh}}(p)\).  Then
\begin{equation}\label{eq:generalSignLaw}
 \operatorname{sgn}\Delta_{p,\beta}(s)
 =\begin{cases}
  1, & 0\leq s<s_*,\\
  0, & s=s_*,\\
  -1, & s_*<s<1.
 \end{cases}
\end{equation}
At \(s=s_*\), the maximizing energy is unique and equals
\begin{equation}\label{eq:generalContactEnergy}
 u^*_{p,\beta}
 =\frac{2(p-1)}{p-2}\,\beta s_*^{p/2}.
\end{equation}
In particular,
\(u^*_{p,\beta_{\mathrm{sh}}(p)}=u_{\infty,p}\).
\end{proposition}

\begin{proof}
Direct integration in \eqref{eq:Omega} gives
\begin{equation}\label{eq:OmegaExplicit}
 \Om(x)=
 \begin{cases}
  \dfrac{x^2}{4}-\dfrac12, & |x|\leq2,\\[6pt]
  \dfrac{x^2}{4}-\dfrac12-
  \dfrac{|x|\sqrt{x^2-4}}{4}
  +\log\!\left(\dfrac{|x|+\sqrt{x^2-4}}2\right), & |x|\geq2.
 \end{cases}
\end{equation}
Consequently, on the quadratic branch,
\begin{equation}\label{eq:generalSigmaInside}
 \Sigma_p(u)
 =\frac12\log(p-1)-\frac{p-2}{4(p-1)}u^2,
 \qquad |u|\leq u_{\infty,p}.
\end{equation}
For \(u=u_{\infty,p}\cosh a\), \(a\geq0\), the same formula gives
\[
 -\Sigma_p'(u)
 =\frac{(p-2)\cosh a+p\sinh a}{\sqrt{p(p-1)}}.
\]
Together with \eqref{eq:generalSigmaInside}, this shows that
\(-\Sigma_p'\) is strictly increasing on \([0,\infty)\), with matching
one-sided values at \(u_{\infty,p}\).  By evenness, \(\Sigma_p\) is
strictly concave; it is also coercive by \eqref{eq:OmegaExplicit}.
Hence, after adding any linear mark, the maximizer is unique.
Since the coefficient of \(u\) in \(\Phi_{p,\beta}(u,\sqrt s)\) is
nonnegative, the unique maximizer is nonnegative.  Whenever it lies in
the quadratic branch, it is
\begin{equation}\label{eq:generalInnerOptimizer}
 u_{p,\beta}(s)
 =\frac{2(p-1)}{p-2}\,\beta s^{p/2},
\end{equation}
and substitution in \eqref{eq:PhiGeneral} gives
\begin{equation}\label{eq:general-p-inner}
 \Delta_{p,\beta}(s)
 =\frac12\log\bigl((p-1)(1-s)\bigr)
 +\frac{p\beta^2}{2}s^{p-1}
  \left(\frac{p-1}{p-2}s-1\right).
\end{equation}
The identity
\begin{equation}\label{eq:optimizerEdgeRatio}
 \frac{u_{p,\beta}(s)}{u_{\infty,p}}
 =\frac{\beta}{\beta_{\mathrm{sh}}(p)}
  \left(\frac{s}{s_*}\right)^{p/2}
\end{equation}
shows that \eqref{eq:general-p-inner} applies throughout
\(0\leq s\leq s_*\).  It immediately gives
\(\Delta_{p,\beta}(s_*)=0\) and \eqref{eq:generalContactEnergy}.

We next prove positivity for \(s<s_*\).  Put
\[
 n=p-2,
 \qquad
 x=\frac{s}{s_*}.
\]
At \(\beta=\beta_{\mathrm{sh}}(p)\), equation
\eqref{eq:general-p-inner} becomes
\begin{equation}\label{eq:generalInnerNormalized}
 2\Delta_{p,\beta_{\mathrm{sh}}(p)}(s)
 =\log\bigl(1+n(1-x)\bigr)-nx^{n+1}(1-x).
\end{equation}
For \(t=1-x\in(0,1)\),
\[
 \log(1+nt)>\frac{nt}{1+nt}>nt(1-t)^{n+1}.
\]
The second inequality follows from
\((1-t)^{-(n+1)}>1+nt\).  Thus the right-hand side of
\eqref{eq:generalInnerNormalized} is positive; the case \(t=1\) is
immediate from \(\log(1+n)>0\).  For \(s<s_*\), the
coefficient of \(\beta^2\) in \eqref{eq:general-p-inner} is negative,
and hence
\[
 \Delta_{p,\beta}(s)
 \geq\Delta_{p,\beta_{\mathrm{sh}}(p)}(s)>0.
\]

It remains to treat \(s>s_*\).  We first work at
\(\beta=\beta_{\mathrm{sh}}(p)\).  By
\eqref{eq:optimizerEdgeRatio}, the formal quadratic stationary point
then lies strictly beyond \(u_{\infty,p}\).  Strict concavity therefore
places the actual maximizer in the outer branch.  Write \(s=s_*x\) and
parametrize it by
\[
 u=u_{\infty,p}\cosh a,
 \qquad
 y=e^a.
\]
Formula \eqref{eq:OmegaExplicit} yields
\begin{equation}\label{eq:generalSigmaOutside}
 \Sigma_p(u_{\infty,p}\cosh a)
 =\frac12\log(p-1)-\frac{p-2}{2p}+a
  -\frac{p-1}{2p}e^{2a}+\frac1{2p}e^{-2a}.
\end{equation}
The critical-point equation for the maximizing \(a\) becomes
\begin{equation}\label{eq:generalOuterCritical}
 (n+1)y-y^{-1}=nx^{(n+2)/2}.
\end{equation}
Set
\[
 L=n+1-nx,
 \qquad
 z=x^{n/2}.
\]
Since \(1<x<(n+1)/n\), we have \(L>0\).  Envelope differentiation,
followed by substitution of \eqref{eq:generalOuterCritical}, gives
\begin{equation}\label{eq:general-p-outer-derivative}
 \frac{\dd}{\dd x}
 \Delta_{p,\beta_{\mathrm{sh}}(p)}(s_*x)
 =\frac{n}{2L}(y-zL)(zL-y^{-1}).
\end{equation}
For \(x>1\),
\[
 (zL)'=nz\left(\frac{L}{2x}-1\right)<0,
\]
so \(zL\) decreases from one.  Since \(y>1\), this gives
\(y-zL>0\).  To determine the other factor, define
\[
 f(v)=(n+1)v-v^{-1}-nx^{(n+2)/2}.
\]
The function \(f\) is strictly increasing, \(f(y)=0\), and a direct
calculation gives
\[
 zL\,f\bigl((zL)^{-1}\bigr)
 =(n+1)(1-x^nL)>0,
\]
because
\[
 (x^nL)'=n(n+1)x^{n-1}(1-x)<0,
\]
so \(x^nL\) decreases from one on this interval.  Hence
\(y<(zL)^{-1}\), so \(zL-y^{-1}<0\).  It follows from
\eqref{eq:general-p-outer-derivative} that
\[
 \Delta_{p,\beta_{\mathrm{sh}}(p)}(s)<0,
 \qquad s>s_*.
\]

Finally fix \(s>s_*\) and let \(u_{p,\beta}(s)\) be the maximizing
energy.  The envelope theorem gives
\begin{equation}\label{eq:generalBetaDerivative}
 \partial_\beta\Delta_{p,\beta}(s)
 =s^{p/2}\left[
  u_{p,\beta}(s)
  -\beta s^{(p-2)/2}\bigl(p-(p-1)s\bigr)
 \right].
\end{equation}
This derivative is strictly positive.  In the quadratic branch, the
claim reduces by \eqref{eq:generalInnerOptimizer} to \(s>s_*\).  In
the outer branch, \(u_{p,\beta}(s)\geq u_{\infty,p}\), while
\[
 s\longmapsto s^{(p-2)/2}\bigl(p-(p-1)s\bigr)
\]
is strictly decreasing for \(s>s_*\), and
\[
 2\beta_{\mathrm{sh}}(p)s_*^{(p-2)/2}=u_{\infty,p}.
\]
Since \(\beta\leq\beta_{\mathrm{sh}}(p)\), the bracket in
\eqref{eq:generalBetaDerivative} is again positive.  Therefore
\[
 \Delta_{p,\beta}(s)
 \leq\Delta_{p,\beta_{\mathrm{sh}}(p)}(s)<0,
\]
which completes the proof.
\end{proof}

\end{document}
